%% comparison_section.tex -- DRAFT "Comparison with previous work" subsection for
%% manuscripts/CV_TimeSeries/main.tex (Discussion).  Committee package cv-literature,
%% findings ED.E1 / RF.M8 / CV-R12, with PH.P2, P3, P5, P6 and RF minor 1 folded in.
%%
%% HOW TO USE.  \input this file inside Section "Discussion" of main.tex, after
%% "What these data add", and %% literature_macros.tex -- GENERATED by committee/work/cv-literature/literature_scales.py.
%% Do not edit.  External (literature) constants and arithmetic on them; none is a measurement
%% of this programme.  Each line names its source, in the manner of numbers.tex.
\newcommand{\LitStLmiCropperPeriodD}{0.07908908}  % [literature: cropper1986, Sect. 3]
\newcommand{\LitStLmiCropperPeriodErrD}{8 \times 10^{-8}}  % [literature: cropper1986, Sect. 3]
\newcommand{\LitStLmiCropperEpoch}{2445059.7024}  % [literature: cropper1986, Sect. 3, ephemeris]
\newcommand{\LitStLmiBfield}{12.1}  % [literature: campbell2008c, abstract; ST LMi modelling results]
\newcommand{\LitStLmiVsxMinusCropperSigma}{0.5}  % [derived in literature_scales.py: (P_VSX - P_Cropper) / sigma_Cropper]
\newcommand{\LitStLmiRefitMinusCropperSigma}{0.6}  % [derived in literature_scales.py: (P_refit - P_Cropper) / sqrt(sigma_refit^2 + sigma_Cropper^2)]
\newcommand{\LitStLmiCropperDriftCycles}{0.022}  % [derived in literature_scales.py: NumStLmiCycles x sigma_P(Cropper) / P: phase drift over the paper's cycle count under the only PUBLISHED period uncertainty]
\newcommand{\LitStLmiCropperDriftMargin}{23}  % [derived in literature_scales.py: 0.5 cycle / that drift]
\newcommand{\LitStLmiVsxZeroOnCropper}{0.97}  % [derived in literature_scales.py: fractional part of the above: polarimetric phase of VSX phase zero]
\newcommand{\LitStLmiVsxZeroOnCropperErr}{0.18}  % [derived in literature_scales.py: cycles x sigma_P / P (1 sigma; epoch error and HJD-vs-BJD, <0.01 cycle, neglected)]
\newcommand{\LitStLmiEdgeOnCropper}{0.13}  % [derived in literature_scales.py: VSX zero + NumStLmiOcOffsetCycles, polarimetric phase]
\newcommand{\LitStLmiEdgeOnConj}{0.96}  % [derived in literature_scales.py: VSX zero + NumStLmiOcOffsetCycles, conjunction phase]
\newcommand{\LitStLmiVsxZeroOnConjErr}{0.11}  % [derived in literature_scales.py: cycles x sigma_P / P]
\newcommand{\LitStLmiCropperBaselineYr}{44}  % [derived in literature_scales.py: that baseline in years]
\newcommand{\LitStLmiFallRangeS}{410--478}  % [derived in literature_scales.py: 0.06 P (Peacock et al. 1992) ... (1.05 - 0.98) P (Cropper 1986)]
\newcommand{\LitStLmiSeasonShiftS}{171}  % [derived in literature_scales.py: half the 1982->1985 change in bright-phase duration: how far one edge moves if the change is symmetric]
\newcommand{\LitStLmiSeasonShiftDeg}{9}  % [derived in literature_scales.py: the same in degrees]
\newcommand{\LitStLmiChromWhiteJS}{68}  % [derived in literature_scales.py: half the white-light -> J difference in duration (Bailey et al. 1985)]
\newcommand{\LitStLmiChromJKS}{273}  % [derived in literature_scales.py: half the J -> K difference]
\newcommand{\LitStLmiChromGIS}{37}  % [derived in literature_scales.py: the white->J edge shift scaled by ln(lambda_i/lambda_g)/ln(1.25/0.55): a log-linear interpolation, ORDER OF MAGNITUDE ONLY; sign: redder band ends later]
\newcommand{\LitStLmiHarmonicG}{18}  % [derived in literature_scales.py: lambda_c / lambda_g]
\newcommand{\LitStLmiHarmonicR}{14}  % [derived in literature_scales.py: lambda_c / lambda_r]
\newcommand{\LitStLmiHarmonicI}{12}  % [derived in literature_scales.py: lambda_c / lambda_i]
\newcommand{\LitStLmiStateRecurrenceD}{273 \pm 90}  % [literature: duffy2022, Table 1]
\newcommand{\LitStLmiStateDurationD}{24.1 \pm 7.6}  % [literature: duffy2022, Table 1]
\newcommand{\LitVfifteenPdot}{3.86 \times 10^{-8}}  % [literature: schmidt1995, quoted in Pavlenko et al. (2018) Sect. 1 (polarimetry, 1987-1992)]
\newcommand{\LitVfifteenOverBound}{11}  % [derived in literature_scales.py: Pdot_spin(V1500 Cyg) / NumStLmiPdotLimit]
\newcommand{\LitPdotSecularRange}{4 \times 10^{-14}$--$2 \times 10^{-13}}  % [derived in literature_scales.py: P / (5 Gyr): secular orbital evolution, slow end]
\newcommand{\LitLibrationPdot}{4 \times 10^{-11}}  % [derived in literature_scales.py: peak curvature of a DP Leo-like libration (A=25 deg, T=60 yr) expressed as an apparent Pdot at ST LMi's period: A (2 pi/T)^2 / 360 x P^2]
\newcommand{\LitBoundOverLibration}{101}  % [derived in literature_scales.py: NumStLmiPdotLimit / that]
\newcommand{\LitDpLeoRateDegYr}{2.3}  % [derived in literature_scales.py: 50 deg / (2001 - 1979)]
\newcommand{\LitStLmiRmsDeg}{4.4}  % [derived in literature_scales.py: NumStLmiOcRmsS / P x 360: O-C scatter as spot longitude]
\newcommand{\LitYzVpPeriodD}{0.086924}  % [literature: vanparadijs1994, Sect. 6]
\newcommand{\LitYzVpPeriodErrD}{0.000007}  % [literature: vanparadijs1994, Sect. 6]
\newcommand{\LitYzDayDriftSh}{0.027}  % [derived in literature_scales.py: phase drift over 1 d under Shafter & Hessman's published sigma_P (cycles)]
\newcommand{\LitYzDayDriftVp}{0.0009}  % [derived in literature_scales.py: the same under van Paradijs et al.'s sigma_P]
\newcommand{\LitYzDaySlip}{0.016}  % [derived in literature_scales.py: phase slip per day from folding on 0.0868 d instead of 0.086924 d]
\newcommand{\LitYzSuperAmp}{4.0}  % [derived in literature_scales.py: 15.0 - 11.0 (Kato et al. 2002)]
\newcommand{\LitYzNormalAmp}{3.0}  % [derived in literature_scales.py: 15.2 - 12.2 (Hakala et al. 2004)]
\newcommand{\LitYzHumpSemiVpMmag}{250}  % [derived in literature_scales.py: half of 0.5 mag, in mmag]
\newcommand{\LitYzDeclineSemiVpMmag}{100}  % [derived in literature_scales.py: half of 0.2 mag, in mmag]
\newcommand{\LitShSemiPeakMmag}{125}  % [derived in literature_scales.py: half of 0.25 mag, in mmag]
\newcommand{\LitYzShSemiTessMmag}{150}  % [derived in literature_scales.py: half of 0.3 mag, in mmag]
\newcommand{\LitYzPshKato}{0.09031(5)}  % [literature: kato2009, Sect. 6.34]
\newcommand{\LitYzPshTess}{0.09043(27)}  % [literature: dai2026, abstract]
\newcommand{\LitYzSupercycleD}{134 \pm 19}  % [literature: patterson1979, Table I]
\newcommand{\LitStLmiFaintFraction}{0.64}  % [literature: cropper1986, Sect. 4]
\newcommand{\LitStLmiFallPhases}{0.98--1.05}  % [literature: cropper1986, Sect. 4]
\newcommand{\LitStLmiFallEndPhase}{0.05}  % [literature: cropper1986, Sect. 4]
\newcommand{\LitStLmiFallWidthPeacock}{0.06}  % [literature: peacock1992, results, R_c light curve of 1986 Jan 8]
\newcommand{\LitStLmiFallWidthRange}{0.06--0.07}  % [literature: peacock1992, results, R_c light curve of 1986 Jan 8; literature: cropper1986, Sect. 4]
\newcommand{\LitStLmiDropPhaseRob}{0.87}  % [literature: robertson2008, photometry, 2007 high state]
\newcommand{\LitStLmiBfieldFerrario}{12}  % [literature: ferrario1993, quoted in Kafka et al. (2007) Sect. 1 (+-0.5 MG); '~12.0 MG' in Campbell et al. (2008)]
\newcommand{\LitStLmiCycFundamentalUm}{9}  % [derived in literature_scales.py: 1.0710 um x (100 MG / 12.1 MG)]
\newcommand{\LitStLmiSfstPeriod}{0.079087 \pm 0.000001}  % [literature: stockman1983, quoted in Cropper (1986) Sect. 3]
\newcommand{\LitStLmiBwsPeriodD}{0.0790898}  % [literature: bailey1985, Sect. 2.3]
\newcommand{\LitStLmiDurCropperEarly}{0.33}  % [literature: cropper1986, Sect. 4]
\newcommand{\LitStLmiDurCropperLate}{0.38}  % [literature: cropper1986, Sect. 4]
\newcommand{\LitStLmiDurBaileyEarly}{0.24}  % [literature: bailey1985, Sect. 8]
\newcommand{\LitStLmiDurBaileyLate}{0.29}  % [literature: bailey1985, Sect. 8]
\newcommand{\LitStLmiDurWhite}{0.29}  % [literature: bailey1985, Sect. 7.3]
\newcommand{\LitStLmiDurJ}{0.31}  % [literature: bailey1985, Sect. 7.3]
\newcommand{\LitStLmiDurH}{0.36}  % [literature: bailey1985, Sect. 7.3]
\newcommand{\LitStLmiDurK}{0.39}  % [literature: bailey1985, Sect. 7.3]
\newcommand{\LitStLmiExtentDeg}{15--18}  % [literature: potter2000, Stokes-imaging solution]
\newcommand{\LitStLmiLowV}{17.5}  % [literature: kafka2007, Sect. 1 (RoboScope, Kafka & Honeycutt 2005)]
\newcommand{\LitStLmiExtremeLowV}{18.5}  % [literature: kafka2007, abstract]
\newcommand{\LitDpLeoAmpDeg}{25}  % [literature: beuermann2014, abstract]
\newcommand{\LitDpLeoPeriodYr}{60}  % [literature: beuermann2014, abstract]
\newcommand{\LitYzTnTess}{8--9}  % [literature: sun2026, Sect. 3.16]
\newcommand{\LitYzTnPatterson}{10.3 \pm 2.5}  % [literature: patterson1979, Table I]
\newcommand{\LitYzKatoMax}{11.0}  % [literature: kato2002, Table 3]
\newcommand{\LitYzKatoMin}{15.0}  % [literature: kato2002, Table 3]
\newcommand{\LitShAmpMaxMag}{0.25}  % [literature: smak2010, abstract]
\newcommand{\LitYzShAmpTessMag}{0.3}  % [literature: dai2026, Sect. 4.2]
\newcommand{\LitYzHumpFullVpMag}{0.5}  % [literature: vanparadijs1994, Sect. 6]
\newcommand{\LitYzDeclineFullVpMag}{0.2}  % [literature: vanparadijs1994, Sect. 7]
\newcommand{\LitYzHumpMaxPhase}{0.8}  % [literature: vanparadijs1994, Sect. 6]
\newcommand{\LitYzDeclineMaxPhase}{0.5}  % [literature: vanparadijs1994, Sect. 7]
\newcommand{\LitYzEpsPct}{4.0}  % [literature: kato2009, Sect. 6.34]
\newcommand{\LitYzShPeriod}{0.0868(2)}  % [literature: shafter1988, VSX OID 4857 (from Downes et al. 2001); value of Shafter & Hessman (1988) as quoted by Verbunt et al. (1999) and Dai et al. (2026)]
\newcommand{\LitYzVpPeriodParen}{0.086924(7)}  % [literature: vanparadijs1994, Sect. 6]
\newcommand{\LitYzFlickerPp}{0.75}  % [literature: moffett1974, quoted in van Paradijs et al. (1994) Sect. 6 and Zhao et al. (2005)]
\newcommand{\LitStLmiBfieldErr}{0.5}  % [literature: campbell2008c, abstract; ST LMi modelling results]
\newcommand{\LitStLmiIrHarmonics}{4--7}  % [literature: campbell2008c, ST LMi modelling results (n = 4-7 at 2.25-1.30 um)]
\newcommand{\LitStLmiVsxEpoch}{2459298.4236}  % [literature: watson2006, VSX OID 17253 (VizieR B/vsx; vsx.aavso.org), rev. 5 of 2025-11-12 by S. Otero: 'Period and epoch from AAVSO data']
\newcommand{\LitVvPupEpoch}{2427889.6474}  % [literature: walker1965, VSX OID 26642 remark 'The epoch of light maximum is given'; ephemeris reproduced in Howell et al. (2006) Sect. 3]
\newcommand{\LitVvPupPeriodD}{0.0697468256}  % [literature: walker1965, same]
\newcommand{\LitAnUmaVsxEpoch}{2456725.4053}  % [literature: watson2006, VSX OID 37171, rev. 4 of 2025-11-12: 'Epoch from AAVSO data'; remark: 'The epoch of Min ... is given']
\newcommand{\LitAnUmaOkPeriodD}{0.079752867(12)}  % [literature: ok2025, Eq. 1]
\newcommand{\LitEuUmaVsxEpoch}{2458231.7799}  % [literature: chen2020, VSX OID 37268, rev. 5 of 2021-10-26: 'Updated from Chen 2020 catalog']
\newcommand{\LitYzVpEpoch}{2447518.7255}  % [literature: vanparadijs1994, Sect. 6]
\newcommand{\LitYzVpEpochDriftCycles}{12}  % [derived in literature_scales.py: cycles of phase uncertainty in propagating the 1988 epoch to 2024: absolute phase is lost]
 in the preamble beside
%% \input{numbers}.  This package does not edit main.tex; the chair places it.
%%
%% WHERE THE NUMBERS COME FROM.  Three kinds appear and they are kept apart:
%%   \Num...   our measurements, from numbers.tex (the products databases);
%%   \Lit...   literature constants and arithmetic on them, from literature_macros.tex,
%%             emitted by committee/work/cv-literature/literature_scales.py -- each
%%             macro's trailing comment names the paper and the place it was read;
%%   \PH{...}  a value of OURS that does not exist as a macro yet, or that the
%%             revision packages (CV-R1...R6) are about to change.  Each placeholder
%%             names the package that owns it.
%% NO number in this file is typed by hand (calendar years excepted).
%%
%% Every \cite key below exists in references.bib and was verified against its
%% registry of record (bib_verification.md).
\providecommand{\PH}[1]{\textbf{[NUMBER: #1]}}

\subsection{Comparison with previous work} \label{sec:comparison}

\subsubsection{ST~LMi}

\paragraph{The bright phase and the feature we time.}
ST~LMi was identified as an AM~Her system by \citet{stockman1983} and
modelled as a single accreting pole that the white dwarf's limb hides for
most of the orbit \citep{schmidt1983, bailey1985, cropper1986}.  Our folded
profile reproduces that literature closely.  \citet{cropper1986} found the
star faint for \LitStLmiFaintFraction{} of the orbit, a rise to maximum
more gradual than the fall, and a rapid decline to the faint level
occupying polarimetric phases \LitStLmiFallPhases{};
\citet{peacock1992} measured the same fall as taking
\LitStLmiFallWidthPeacock{} of the period in $R_{\rm c}$; and
\citet{robertson2008} again saw a gradual brightening and a quicker fade in
the 2007 high state.  We measure the steepest faintward gradient at
\NumStLmiFallRiseRatioRange{} times the steepest brightward one in every
solved $g/r/i$ series of both eras (Section~\ref{sec:stlmi}), so the
asymmetry is a persistent property of the system and not of one epoch.
The published widths of the fall, \LitStLmiFallWidthRange{} of the orbit,
are \LitStLmiFallRangeS{}~s: the feature we time is a ramp several hundred
seconds long, sampled by exposures whose half-length reaches
\NumPolarHalfExposureS{}~s, which is the physical reason an analytic
single-cycle timing budget is optimistic (Section~\ref{sec:timing}).

The identification of the feature matters for the constant offset of
Section~\ref{sec:oc}.  The polarimetric ephemeris of \citet{cropper1986},
${\rm HJD}\;\LitStLmiCropperEpoch + \LitStLmiCropperPeriodD\,E$, places
phase zero at the linear-polarisation pulse, which that paper locates half
way through the rapid decline --- that is, at the feature we time ---
and \citet{bailey1985} tied their period to the end of the bright phase.
The catalogue ephemeris we fold on is different in kind: its period and
epoch were revised in VSX from AAVSO photometry \citep{watson2006}, carry
no published uncertainty, and the fiducial the epoch marks is not
documented.  Carrying the catalogue epoch back to Cropper's with his period
puts catalogue phase zero at polarimetric phase
$\LitStLmiVsxZeroOnCropper \pm \LitStLmiVsxZeroOnCropperErr$, and our timed
edge, \NumStLmiOcOffsetCycles{} of a cycle later, at
$\LitStLmiEdgeOnCropper \pm \LitStLmiVsxZeroOnCropperErr$ --- consistent
with the pulse at phase zero and the end of the decline at
\LitStLmiFallEndPhase{}, where the uncertainty is Cropper's period error
accumulated over \LitStLmiCropperBaselineYr{}~yr and nothing of ours.  On
the inferior-conjunction ephemeris used by \citet{kafka2007} and
\citet{robertson2008} the same edge falls at phase
$\LitStLmiEdgeOnConj \pm \LitStLmiVsxZeroOnConjErr$, against the sharp
drop those authors report near \LitStLmiDropPhaseRob{}.  The
\NumStLmiOcOffsetS{}~s offset of Section~\ref{sec:oc} is therefore what the
literature leads one to expect if the catalogue epoch marks the
bright-phase maximum and our edge the end of the bright phase; we state it
as a consistency, not as an identification of the catalogue's fiducial,
which only its compiler can supply.

\paragraph{Colour through the orbit.}
That the bright phase of ST~LMi is red is among the oldest facts about it:
the orbital modulation is strongest at red wavelengths
\citep{stockman1983, cropper1986}, larger in $R_{\rm c}$ and $I_{\rm c}$
\citep{peacock1992} than in $V$ \citep{kafka2007}, and \citet{kafka2007}
note explicitly that the star reddens at the phase of the bright hump
because the cyclotron continuum then outshines every other component.  The
cause is the field strength.  Infrared cyclotron humps give
$B \simeq \LitStLmiBfieldFerrario$~MG at the accreting pole
\citep{ferrario1993}, refined to
$\LitStLmiBfield \pm \LitStLmiBfieldErr$~MG by phase-resolved modelling
\citep{campbell2008c}; the fundamental then lies near
\LitStLmiCycFundamentalUm~$\mu$m, harmonics \LitStLmiIrHarmonics{} fall in
the near infrared where \citet{campbell2008c} resolve them, and our $i$,
$r$ and $g$ bands sample harmonics near \LitStLmiHarmonicI{},
\LitStLmiHarmonicR{} and \LitStLmiHarmonicG{}.  At such high harmonics the
emission is optically thin and its intensity falls steeply with harmonic
number \citep{chanmugam1981, wickramasinghe1985, wickramasinghe2000}, so
the cyclotron component is far stronger in $i$ than in $g$ and the system
reddens whenever the accretion region is in view.
Figure~\ref{fig:colourphase} is the state-tagged, many-night version of
that statement: $g-r$ swings redward by
\PH{g-r bright-minus-faint amplitude, mag, per era; CV-R6}
through the bright phase, with a turn-over phase that repeats between the
two eras to
\PH{two-era repeatability of the colour-extremum phase, cycles; CV-R6},
and the swing survives tightening the pairing window from
\NumColourPairWindowS{}~s to \PH{result with a <=120 s pairing window; CV-R6}.
We are careful about what is new here.  Simultaneous multiband light curves
of single nights exist --- \citet{bailey1985} in the optical and near
infrared, \citet{peacock1992} with polarimetry on three occasions --- and
the sign of the colour change has been known for four decades.  What this
programme adds is the number of nights
(\NumStLmiThreeFilterNights{} three-filter full-orbit nights), their
spread over two seasons, and an accretion state for each night taken from
an independent survey record.

\paragraph{The period, and what the $\dot{P}$ bound constrains.}
Three published periods precede ours: $\LitStLmiSfstPeriod$~d from
polarisation peaks in 1982 \citep{stockman1983}, \LitStLmiBwsPeriodD~d
tied to the end of the bright phase \citep{bailey1985}, and
$\LitStLmiCropperPeriodD \pm \LitStLmiCropperPeriodErrD$~d from
polarisation peaks spanning 1982--1985 \citep{cropper1986}, the value
adopted by later studies \citep[e.g.][]{kafka2007, robertson2008}.  The
catalogue period we fold on differs from Cropper's by
\LitStLmiVsxMinusCropperSigma{} of his standard errors, and our own refit,
$\NumStLmiFittedPeriodD \pm \NumStLmiFittedPeriodSigmaD$~d, by
\LitStLmiRefitMinusCropperSigma$\sigma$ of the combined error.  This also
settles a point our cycle count left open.  Section~\ref{sec:oc} argued
uniqueness from an \emph{assumed} period error, the last digit of the
catalogue string, because the catalogue publishes none.  The only published
error is Cropper's, and under it the phase drift over the
\NumStLmiCycles{} cycles to our last edge is \LitStLmiCropperDriftCycles{}
cycles: the count is unique by a factor of \LitStLmiCropperDriftMargin{},
not by the larger margin the assumed error gave.

Every one of these periods was measured from the accretion region ---
polarisation pulses, or the bright-phase edge --- so each is the period on
which that region returns to the limb: the white dwarf's rotation, plus any
drift of the region in magnetic longitude.  In a synchronised polar that
equals the orbital period, but it is not a measurement of the orbit, and
our bound should be read accordingly: $|\dot{P}| < \NumStLmiPdotLimit$ is a
limit on the \emph{spin and spot-longitude} period of ST~LMi over
\NumStLmiEpochSpanYears{}~yr.  Three physical scales say what such a limit
can and cannot reach.  (i) Secular orbital evolution below the period gap
proceeds on gigayear timescales \citep{knigge2011}, i.e.\
$|\dot{P}| \sim \LitPdotSecularRange$, four to five orders of magnitude
below the bound; nothing here tests it.  (ii) A white dwarf knocked out of
synchronism relaxes far faster: V1500~Cyg, desynchronised by its 1975 nova
\citep{stockman1988, schmidt1991}, showed
$\dot{P}_{\rm spin} = \LitVfifteenPdot$ in 1987--1992 \citep{schmidt1995},
with resynchronisation expected on a timescale of order a century or two
\citep{schmidt1995, campbell1999} and its present state still debated
\citep{harrison2016, harrison2018, pavlenko2018}; CD~Ind and V1432~Aql are
the other well-studied cases \citep{myers2017, staubert2003}.  Our bound is
\LitVfifteenOverBound{} times smaller than the V1500~Cyg rate, so ST~LMi is
not relaxing in that way.  (iii) A synchronised white dwarf can still
librate about its equilibrium orientation: in DP~Leo the accretion spot
advanced in longitude by about \LitDpLeoRateDegYr$^{\circ}$~yr$^{-1}$ for
two decades and then stopped, consistent with an oscillation of
$\sim$\LitDpLeoAmpDeg$^{\circ}$ amplitude and $\sim$\LitDpLeoPeriodYr~yr
period \citep{beuermann2014}.  A steady drift is absorbed by the fitted
period and is invisible to a $\dot{P}$ test; only its curvature is not, and
a libration of DP~Leo's size would appear at ST~LMi's period as
$|\dot{P}| \lesssim \LitLibrationPdot$, a factor of
\LitBoundOverLibration{} below our bound.  The bound therefore excludes
scale (ii) and constrains neither (i) nor (iii).

The statement these data do support is about longitude.  The
\NumStLmiOcRmsS{}~s scatter of the published epochs is
\LitStLmiRmsDeg$^{\circ}$ of white-dwarf rotation, and that is small
against what the literature says the edge can do: the bright phase
lengthened from \LitStLmiDurCropperEarly{} to \LitStLmiDurCropperLate{} of
the orbit between 1982 and 1985 \citep{cropper1986}, and from
\LitStLmiDurBaileyEarly{} to \LitStLmiDurBaileyLate{} within a single year
\citep{bailey1985}, which moves one edge by about \LitStLmiSeasonShiftS{}~s
(\LitStLmiSeasonShiftDeg$^{\circ}$) if the change is symmetric.
\PH{One sentence, CV-R3/CV-R5: the rms of the night-level epochs in degrees (N = 17 nights, era-offset nuisance fitted), the edge-longitude difference between high- and low-state nights with its error, and whether either reaches the 9-degree scale; then the sensitivity of both to dropping the 2024 epochs.}

\paragraph{Should the edge epoch depend on wavelength?}
It should, and the literature gives the sign and the scale.
\citet{bailey1985} found the bright phase to last \LitStLmiDurWhite{} of
the orbit in white light but \LitStLmiDurJ{}, \LitStLmiDurH{} and
\LitStLmiDurK{} in $J$, $H$ and $K$, in simultaneous data, and attributed
the trend to emission from greater heights at longer wavelengths, where the
cyclotron opacity is larger.  A longer bright phase ends later: half the
white-to-$J$ difference is \LitStLmiChromWhiteJS{}~s, half the $J$-to-$K$
difference \LitStLmiChromJKS{}~s.  The accretion region is also extended,
by some $\pm$\LitStLmiExtentDeg$^{\circ}$ of magnetic longitude in Stokes
imaging \citep{potter2000} and with resolved structure in other
reconstructions \citep{cropper1994, stockman1996}, so bands that weight it
differently need not share a half-light point.  A log-linear interpolation
of Bailey et al.'s durations to our $g$ and $i$ bands gives of order
\LitStLmiChromGIS{}~s, with the redder band's edge the later; we quote it
as an order of magnitude only, because the infrared trend is an
optical-depth effect and our bands are optically thin.  That is the
predicted scale our band-pair test should be read against.
\PH{CV-R1: the combined-era g-i edge offset in seconds, with its scatter-based error and sign.}
\PH{CV-R2: the signed matched-cell injection bias for the same pair.}
\PH{D3: one sentence stating which the injection test supports -- an offset of the sign and order the literature predicts, or an estimator bias of that size.}

\paragraph{Accretion states.}
ST~LMi's long-term behaviour is unusually well documented.  RoboScope
photometry \citep{kafka2005} shows a five-year low state (1992--1997) near
$V = \LitStLmiLowV$ followed by six years of high state; \citet{kafka2007}
add an extreme low state near $V = \LitStLmiExtremeLowV$
\citep[see also][]{howell2000}; the Catalina record shows three brightness
levels \citep{mason2017}; and ZTF resolves short-lived states lasting
$\LitStLmiStateDurationD$~d that recur every $\LitStLmiStateRecurrenceD$~d
\citep{duffy2022}.  Our state classification (Section~\ref{sec:states}) is
made against the same kind of survey record and should be read with those
timescales in mind: a state lasting three to four weeks is resolved by a
season and not by a night.
\PH{cv-stats: how many of our full-orbit nights fall inside a survey-defined short-lived state.}
What the earlier work could not do, for lack of an independent state record
on the night, is ask whether the colour curve changes \emph{shape} between
states; Figure~\ref{fig:stlmifold} is that comparison, subject to the limits
of Section~\ref{sec:tie}.

\subsubsection{YZ~Cnc}

\paragraph{The outburst record and the superhump literature.}
YZ~Cnc is one of the most active SU~UMa stars known.  \citet{patterson1979},
who discovered its superhumps, found normal outbursts every
$\LitYzTnPatterson$~d and supermaxima every $\LitYzSupercycleD$~d, about a
magnitude brighter than normal maxima and several times longer
\citep[see also][]{szkody1984, vanparadijs1994}; \emph{TESS} recorded
normal outbursts every \LitYzTnTess{}~d \citep{sun2026}.  Catalogued
extremes of \LitYzKatoMin{} and \LitYzKatoMax{}~mag give a superoutburst
amplitude of \LitYzSuperAmp{}~mag \citep{kato2002}, against about
\LitYzNormalAmp{}~mag for a well-observed normal outburst
\citep{hakala2004, vanparadijs1994}.  Our brightest dense run reaches
\NumOutburstPeakAmp{}~mag above quiescence, short of either, which is the
quantitative content of the AAVSO-based classification of
Section~\ref{sec:yzcnc}: the runs lie on the rise and decline of normal
outbursts.  The superhump period itself is well established:
\LitYzPshKato~d in the 2007 superoutburst \citep{kato2009}, which
superseded the discovery value, with later superoutbursts in
\citet{kato2014v, kato2014vi, kato2015}, and \LitYzPshTess~d from two
superoutbursts observed continuously by \emph{TESS}
\citep{dai2026, sun2026}: a fractional excess of \LitYzEpsPct{} per cent
over the orbital period of \citet{vanparadijs1994}, in line with the
excess--period relation of SU~UMa stars \citep{patterson2005} and the
tidal--thermal instability picture behind it
\citep{whitehurst1988, osaki1989, osaki1996, osaki2013}.

\paragraph{What our search could and could not have seen.}
Superhump amplitudes are largest near the start of the plateau and decay
through it \citep{smak2010, kato2012}.  At low inclination the maximum full
amplitude is about \LitShAmpMaxMag{}~mag \citep{smak2010, kato2012}, a
semi-amplitude of \LitShSemiPeakMmag{}~mmag, and \emph{TESS} recorded about
\LitYzShAmpTessMag{}~mag (\LitYzShSemiTessMmag{}~mmag) in YZ~Cnc's own
early plateau \citep{dai2026}; the \NumSuperhumpFloorMmag{}~mmag we adopted
as a floor lies well below both and can describe only a decayed,
late-plateau signal.  A 90 per cent
recovery contour is the \emph{smallest} amplitude a search recovers, so our
blind contours of \NumBlindContourRangeMmag{}~mmag mean this: a freshly
developed superhump of \LitShSemiPeakMmag--\LitYzShSemiTessMmag{}~mmag
would have been recovered in
\PH{CV-R4: the number of run--filters whose contour lies below 125 mmag}
of the \NumBlindContourRunFilters{} run--filter combinations that carry a
contour, and a late-stage superhump near \NumSuperhumpFloorMmag{}~mmag in
none of them.  We therefore exclude large-amplitude superhumps on those
runs and nothing smaller.  The exclusion is also what the literature
predicts.  \citet{dai2026} find no persistent superhump-band signal in
YZ~Cnc's quiescent and normal-outburst intervals; the exception known in
this star is the survival of late superhumps into the normal outburst that
follows a superoutburst \citep{kato2014v}, a behaviour \emph{TESS} has now
shown in three systems \citep{sun2026}.
\PH{From the AAVSO record already held: the interval between the preceding superoutburst and each outburst dense run, and whether any run is the first normal outburst after a superoutburst.}

\paragraph{The orbital hump.}
\citet{vanparadijs1994} reported orbital modulation in YZ~Cnc: a hump of
\LitYzHumpFullVpMag{}~mag full amplitude in quiescence, with maximum near
spectroscopic phase \LitYzHumpMaxPhase{} as expected of a bright spot,
falling to about \LitYzDeclineFullVpMag{}~mag and moving to phase
\LitYzDeclineMaxPhase{} late in the decline from outburst; they folded on
${\rm HJD}\;\LitYzVpEpoch + \LitYzVpPeriodParen\,E$.  Against that
\LitYzHumpSemiVpMmag{}~mmag semi-amplitude our fitted
\NumHumpAmpRangeMmag{}~mmag is smaller by a factor of several and does not
clear the red-noise contour of \NumHumpSelfContourMmag{}~mmag; the Fourier
analysis of the \emph{TESS} superoutburst sectors likewise shows no
orbital signal \citep{sun2026}.  On the
\PH{CV-R4: the number of scopes whose red-noise contour lies below 250 mmag}
of \NumHumpScopesTested{} scopes whose contour lies below
\LitYzHumpSemiVpMmag{}~mmag, a hump as large as that of 1988--89 would have
been recovered and was not, so the quiescent hump of this star is not a
permanent feature.  Two details of our own analysis change in the light of
that paper.  First, a spectroscopic epoch \emph{does} exist, but its period
error accumulates to \LitYzVpEpochDriftCycles{} cycles by 2024, so absolute
orbital phase remains unavailable and only our within-run phase statements
stand.  Second, the period we folded on is the \LitYzShPeriod~d of
\citet{shafter1988}, which is what the catalogues carry
\citep{downes2001, watson2006}; under its published error the phase drift
between consecutive nights is \LitYzDayDriftSh{} cycles, not the smaller
figure the catalogue string implies, and the more precise
\LitYzVpPeriodD~d differs from it by \LitYzDaySlip{} cycles per day.  Both
are far below the \NumHumpPhaseNightShift{}-cycle shift we measure between
the two May nights, so the incoherence between nights is not an artefact of
the ephemeris.
\PH{CV-R4: the between-night phase shift re-derived on the 0.086924 d period.}

\paragraph{Flickering.}
YZ~Cnc's flickering is among the strongest known in cataclysmic variables,
\LitYzFlickerPp{}~mag peak to peak \citep{moffett1974}, and
\citet{vanparadijs1994} attribute the difficulty of earlier hump searches
to it.  Our structure functions \citep[in the sense of][]{simonetti1985}
give \NumFlickerAmpRangeMmag{}~mmag over \NumFlickerTauRangeS{}~s above a
measured floor; \citet{bruch2021} provides the comparative framework
against which such amplitudes can be placed.  It is the reason the relevant
detection threshold for any periodic signal in this star is the red-noise
contour and not the field-star one.

%% --------------------------------------------------------------------------
%% Table: the ephemerides this paper folds on, what each epoch marks, and where
%% each comes from.  Periods are the catalogue values of numbers.tex; epochs
%% and provenance are literature constants (literature_macros.tex).
%% --------------------------------------------------------------------------
\begin{deluxetable*}{llllll}
\tablecaption{Adopted ephemerides: provenance and fiducials.  The catalogue
values are those of the AAVSO International Variable Star Index
\citep[VSX;][]{watson2006} as retrieved for this work; the last column
gives the most precise published ephemeris for comparison.
\label{tab:ephem}}
\tablehead{\colhead{Target} & \colhead{$P$ (d)} & \colhead{$T_0$ (HJD)} &
\colhead{$T_0$ marks} & \colhead{Origin of the catalogue values} &
\colhead{Published ephemeris}}
\startdata
ST~LMi & \NumStLmiPeriodD & \LitStLmiVsxEpoch & not documented &
  VSX, from AAVSO photometry; no error given &
  \citet{cropper1986}: linear-polarisation pulse \\
VV~Pup & \NumVvPupPeriodD & \LitVvPupEpoch & maximum light &
  \citet{walker1965}, as in the General Catalogue \citep{samus2017} &
  \citet{walker1965}; see \citet{howell2006} \\
EU~UMa & \NumEuUmaPeriodD & \LitEuUmaVsxEpoch & minimum light &
  ZTF periodic-variable catalogue \citep{chen2020} &
  \citet{mittaz1992}; \citet{howell1995} \\
AN~UMa & \NumAnUmaPeriodD & \LitAnUmaVsxEpoch & minimum light &
  period: VSX, origin not documented; epoch: VSX, from AAVSO photometry &
  \citet{bonnetbidaud1996}; \citet{ok2025}: \LitAnUmaOkPeriodD~d \\
YZ~Cnc & \NumYzCncPeriodD & \nodata & \nodata &
  \citet{shafter1988} via \citet{downes2001} &
  \citet{vanparadijs1994}: \LitYzVpPeriodParen~d \\
\enddata
\end{deluxetable*}
